% ---------------------------------------------------------------------------
% Informedness of the gradient-informed single-site proposal
%
% Recap of the proposal mathematics, the dt/sqrt(T) scaling law, how dt is
% chosen, and what the dt=460 vs dt=63 pair of Boltzmann-validation runs
% established. Companion to softmax/DEVLOG.txt -- the DEVLOG is the
% chronological record, this is the standing explanation.
%
% Build:  cd softmax/docs && pdflatex informedness.tex && pdflatex informedness.tex
% (twice, for the table of contents and cross-references)
% ---------------------------------------------------------------------------
\documentclass[11pt,a4paper]{article}

\usepackage[T1]{fontenc}
\usepackage[utf8]{inputenc}
\usepackage[margin=2.5cm]{geometry}
\usepackage{amsmath,amssymb}
\usepackage{booktabs}
\usepackage{array}
\usepackage{graphicx}
\usepackage[hidelinks]{hyperref}
% Long file paths are set with hyperref's \path, which breaks at slashes
% and needs no underscore escaping; \sloppy absorbs the rest.
\sloppy

% Consistent notation for the recurring quantities
\newcommand{\dt}{\Delta t}
\newcommand{\gradU}{\nabla U}
\newcommand{\pmatch}{p_{\mathrm{match}}}
\newcommand{\Mass}{M}
\newcommand{\Ncomp}{K_{\mathrm{c}}}

\title{Informedness of the gradient-informed single-site proposal\\
       \large Why it scales as $\dt/\sqrt{T}$, how $\dt$ is chosen,\\
       and what the $\dt=460$ vs.\ $\dt=63$ runs showed}
\author{\texttt{my\_extendedSampling} / \texttt{softmax}}
\date{2026--10--03}

\begin{document}
\maketitle

\begin{abstract}
\noindent
The sampler in \path{classes/rate_sampler.py} proposes one substitution per
Monte Carlo move by reading off which amino acid a gradient-informed Gaussian
displacement ``points at''. How reliably that displacement follows the
gradient's own preference---its \emph{informedness}---is controlled by $\dt$,
but \emph{not} by $\dt$ alone: the governing combination is
$\dt/\sqrt{\Mass T}$. This note derives that scaling from the proposal
construction, records the measured informedness curves for both reference
sequences, explains why $\dt$ is chosen at a deliberately \emph{partial}
informedness of $\pmatch\!\approx\!0.15$ rather than as large as possible, and
reports what the paired $\dt=460$ / $\dt=63$ Boltzmann-validation runs on
\texttt{zero\_polymer} ($L=566$) established. The short version of the last
point: both runs reproduce the exactly-enumerated Boltzmann distribution, so
the $\dt$ error that motivated the pair cost test \emph{coverage}, not
correctness. A subsequent doubling of the run length confirms the residual
discrepancy decays as finite-sample noise must, in two independent
functionals.
\end{abstract}

\tableofcontents

% ===========================================================================
\section{Scope and sources}
% ===========================================================================

Everything below is either derived from the code as it stands or measured in a
run whose output file is committed. Source of truth for each:

\begin{itemize}
  \item Proposal construction: \path{classes/rate_sampler.py}, method
        \texttt{\_step}.
  \item Pointing probability: \path{utils/proposals.py}, function
        \path{log_pointing_prob}.
  \item Informedness curves: \path{tests/test_informedness_vs_dt.py};
        outputs \path{protein_g/informedness_vs_dt_results.txt} and
        \path{zero_polymer/informedness_vs_dt_results.txt}.
  \item Boltzmann validation: \path{tests/validate_boltzmann_toy.py};
        outputs \path{zero_polymer/validate_boltzmann_results_dt460.txt}
        and \path{..._dt63.txt}.
  \item Chronology and the mistake that prompted this note:
        \path{softmax/DEVLOG.txt}, entries dated 2026--10--02 and
        2026--10--03.
\end{itemize}

Quantities that are \emph{estimated} rather than measured are flagged as such
at the point of use; there are two (the effective sample size in
\S\ref{sec:noise}, and the $\pmatch$ interpolations in
Table~\ref{tab:working-points}).

% ===========================================================================
\section{The proposal, and what ``informedness'' means}
% ===========================================================================

\subsection{One move}

A sequence of length $L$ is held as per-site logits over the
$K=25$-letter vocabulary \path{SEQUENCE_USED_VOCAB}, mean-centred along the
class axis to gauge-fix the translation invariance of the softmax. One move:

\begin{enumerate}
  \item Pick one site $s$ uniformly from $\{0,\dots,L-1\}$. Uniform and
        state-independent, so it cancels exactly from the
        Metropolis--Hastings ratio and needs no correction term.
  \item At that site only, draw a displacement $\delta x\in\mathbb{R}^{K}$
        (\S\ref{sec:displacement}) and propose the substitution it points at,
        \emph{excluding} the site's current residue and the five non-canonical
        codes \texttt{X,B,U,Z,O}. That leaves
        \[
          \Ncomp = 25 - 5 - 1 = 19
        \]
        competitors. ($\Ncomp$ is $20$ in the one edge case where the site
        already sits on a non-canonical code, since then there is no current
        residue to remove from the canonical set.)
  \item Accept or reject with an exact Metropolis--Hastings correction using
        the true recomputed energy of the candidate and the true reverse
        proposal probability, so the target is $\exp(-U/T)$ exactly rather
        than to first order.
\end{enumerate}

Note that step 2 always proposes a genuine change: ``staying'' is decided by
rejection in step 3, never by the proposal.

\subsection{The displacement}
\label{sec:displacement}

The displacement is one HMC leapfrog half-step with fresh momentum, which for
a single step is the same object as a MALA/Langevin step:
\begin{equation}
  \delta x \;=\; \frac{\dt\,p}{\Mass} \;-\; \frac{\dt^{2}}{2\Mass}\,\gradU,
  \qquad p \sim \mathcal{N}\!\left(0,\, \Mass T\, I\right),
  \label{eq:leapfrog}
\end{equation}
with $\gradU = \partial U/\partial(\text{logits at } s)$ obtained by relaxing
\emph{only} site $s$ through $\mathrm{softmax}(\cdot/T_{\mathrm{sftm}})$ while
every other site is held at its exact one-hot value. Equation
\eqref{eq:leapfrog} makes $\delta x$ Gaussian, and this is the pair of
quantities everything downstream depends on:
\begin{equation}
  \delta x \sim \mathcal{N}\!\left(\mu,\, \sigma^{2} I\right),
  \qquad
  \boxed{\;\mu \;=\; -\frac{\dt^{2}}{2\Mass}\,\gradU,
         \qquad
         \sigma \;=\; \dt\sqrt{\frac{T}{\Mass}}\;}
  \label{eq:musigma}
\end{equation}
So the drift grows as $\dt^{2}$ while the noise grows only as $\dt$. That
asymmetry is the entire mechanism by which $\dt$ buys informedness, and
\S\ref{sec:scaling} is just this observation carried through carefully.

\subsection{Informedness, defined}

Write $j^{\star}$ for the competitor the gradient alone would pick,
\begin{equation}
  j^{\star} \;=\; \arg\max_{j \in \text{competitors}} \mu_{j}
            \;=\; \arg\min_{j \in \text{competitors}} (\gradU)_{j},
  \label{eq:jstar}
\end{equation}
the second equality because $\mu = -\dt^{2}/(2\Mass)\cdot\gradU$ and
$\dt^{2}/(2\Mass)$ is a \emph{positive scalar}. This gives the first of two
facts that the diagnostic
\texttt{test\_informedness\_vs\_dt.py} was built to separate:

\begin{description}
  \item[(A) The gradient's choice is independent of $\dt$.] Scaling $\mu$ by a
        positive scalar cannot change its $\arg\max$. So \emph{which} residue
        the gradient prefers, and therefore how good that preference is, does
        not depend on $\dt$ at all. $\dt$ cannot improve the gradient's taste.
  \item[(B) How often the realised draw obeys that choice does depend on
        $\dt$.] The actual proposal is the $\arg\max$ of $\delta x$, which is
        $\mu$ plus noise. Define
        \begin{equation}
          \pmatch \;=\; \mathbb{P}\!\left(
            \arg\max_{j} \delta x_{j} = j^{\star} \right),
          \label{eq:pmatch}
        \end{equation}
        the probability that one realised move follows the gradient's own
        pick. This is what $\dt$ controls, and it is what ``informedness''
        means throughout this note.
\end{description}

Conflating (A) and (B) was the original error that this decomposition fixed:
an earlier diagnostic swept $\dt$ looking for improvements in proposal
\emph{quality}, which by (A) cannot exist.

\subsection{The mathematics of (B)}
\label{sec:Bmath}

Statement (B) is not qualitative. $\pmatch$ has a closed form, and every claim
made about it below is a property of that form.

\paragraph{$\pmatch$ is the pointing probability evaluated at the gradient's
own choice.} Comparing \eqref{eq:pmatch} with the definition of $a_{i\to j}$
in \S\ref{sec:aij}, they are the same object with $j$ set to $j^{\star}$:
\begin{equation}
  \pmatch \;=\; a_{i\to j^{\star}}
  \;=\; \int_{-\infty}^{\infty}\! \phi(u)
        \prod_{k \neq i,\,j^{\star}}
        \Phi\!\left(\alpha_{kj^{\star}} + u\right)\mathrm{d}u ,
\end{equation}
using $\beta \equiv 1$ for this sampler (\S\ref{sec:aij}). Substituting
$\mu = -\frac{\dt^{2}}{2\Mass}\gradU$ and $\sigma = \dt\sqrt{T/\Mass}$ makes
the $\dt$ dependence explicit:
\begin{equation}
  \boxed{\;
  \pmatch(\dt)
  \;=\; \int_{-\infty}^{\infty}\! \phi(u)
        \prod_{k \neq i,\,j^{\star}}
        \Phi\!\left(\frac{\dt\, g_{kj^{\star}}}{2\sqrt{\Mass T}} + u\right)
        \mathrm{d}u ,
  \qquad
  g_{kj^{\star}} = (\gradU)_{k} - (\gradU)_{j^{\star}} \;\ge\; 0 ,
  \;}
  \label{eq:pmatch-closed}
\end{equation}
where $g_{kj^{\star}} \ge 0$ holds \emph{by definition} of $j^{\star}$ as the
minimiser of $(\gradU)_{j}$ over the competitors. Four consequences follow,
all verified numerically against \eqref{eq:pmatch-closed}.

\paragraph{1. Monotonicity.} Every $\alpha_{kj^{\star}} = \dt
g_{kj^{\star}}/(2\sqrt{\Mass T})$ is non-decreasing in $\dt$ because
$g_{kj^{\star}}\ge0$; $\Phi$ is increasing; so each factor is non-decreasing
in $\dt$ pointwise in $u$, the product of non-negative non-decreasing
functions is non-decreasing, and integrating against $\phi \ge 0$ preserves
it. Hence $\pmatch$ is \emph{monotonically non-decreasing in} $\dt$ --- there
is no interior optimum to find, which is why choosing $\dt$ is a matter of
picking a working point (\S5) rather than maximising anything. Checked over
$\dt \in [10^{-2},10^{6}]$: monotone at every step.

\paragraph{2. The small-$\dt$ limit is uniform, and equals $1/\Ncomp$
exactly.} As $\dt\to0$ all $\alpha \to 0$ and
\begin{equation}
  \pmatch \longrightarrow \int \phi(u)\,\Phi(u)^{\Ncomp-1}\mathrm{d}u
  \;=\; \frac{1}{\Ncomp} \;=\; \frac{1}{19} \approx 0.0526 ,
\end{equation}
by the standard identity $\int\phi\,\Phi^{n} = 1/(n+1)$ --- the probability
that a designated one of $n+1$ i.i.d.\ variables is the largest. Reproduced to
within quadrature error ($1.4\times10^{-13}$, 120-node Gauss--Hermite). Note this limit is \emph{uniform}, not ``stay put'':
nothing in \eqref{eq:leapfrog} references the current residue's identity,
which is why an earlier design that tried to make small $\dt$ mean ``few sites
change'' could not work and was abandoned.

\paragraph{3. The approach to that floor is linear in $\dt$.} Expanding
$\Phi(\alpha+u) = \Phi(u) + \alpha\phi(u) + O(\alpha^{2})$ in
\eqref{eq:pmatch-closed},
\begin{equation}
  \pmatch(\dt) \;=\; \frac{1}{\Ncomp}
  \;+\; C_{\Ncomp}\,\frac{\dt}{2\sqrt{\Mass T}}
        \sum_{k \neq i,j^{\star}} g_{kj^{\star}}
  \;+\; O(\dt^{2}),
  \qquad
  C_{\Ncomp} = \int \phi(u)^{2}\,\Phi(u)^{\Ncomp-2}\mathrm{d}u ,
\end{equation}
with $C_{19} = 5.393\times10^{-3}$. So near the floor the useful signal grows
linearly in $\dt$, with a slope set by the \emph{sum} of the gradient gaps ---
which is why a sequence with larger gaps reaches a given $\pmatch$ at a
proportionally smaller $\dt$, the effect measured in \S5.3.

Verified against the exact integral on a synthetic 18-competitor case: the
relative error of the linear form falls as the \emph{square} of the largest
shift $\alpha_{\max} = \dt\,g_{\max}/(2\sqrt{\Mass T})$ --- $2.1\times10^{-7}$
at $\alpha_{\max}=1.4\times10^{-3}$, $2.1\times10^{-5}$ at
$1.4\times10^{-2}$, $1.8\times10^{-3}$ at $0.14$, i.e.\ a factor $\sim\!100$
per decade --- which is exactly the $O(\alpha^{2})$ truncation a first-order
expansion should leave. Above $\alpha_{\max}\sim1$ the linear form crosses the
exact curve and then diverges from it, so this is a statement about the floor
only, not a usable approximation at the working point.

\paragraph{4. The large-$\dt$ limit is 1, with one exception.} As
$\dt\to\infty$ every $\alpha_{kj^{\star}} \to +\infty$ provided
$g_{kj^{\star}}>0$, so every factor $\to 1$ and $\pmatch \to 1$: a
deterministic, greedy gradient-descent pick. Verified to machine precision at
$\dt=10^{8}$. The exception is an exact tie, $g_{kj^{\star}}=0$ for some $k$:
that factor stays $\Phi(u)$ for all $\dt$ and the limit becomes $1/2$ (one
tie), or $1/(m+1)$ for $m$ tied competitors --- confirmed numerically at
$0.50000000$ for $m=1$ and $0.33333333$ for $m=2$. This is not a pathology but the correct behaviour, and
it is the formal version of the near-degeneracy that makes a specific reported
substitution unstable between runs (see \path{docs/gradient_method.tex} §6.3).

\paragraph{What $\dt$ therefore buys.} Between those limits, $\dt$ moves the
proposal along a one-parameter family from ``uniform over the 19 competitors''
to ``always the gradient's pick'', monotonically, with no interior structure.
Since by (A) the \emph{identity} of that pick is fixed, $\dt$ selects only how
strictly the sampler commits to it --- and \S5 argues that committing fully is
undesirable, so the right $\dt$ is an interior working point chosen against
the measured curve, not an extremum.

% ===========================================================================
\section{\texorpdfstring{The pointing probability $a_{i\to j}$}{The pointing probability a(i->j)}}
\label{sec:aij}
% ===========================================================================

\subsection{Closed form}

$\pmatch$, and more generally the proposal probability needed for the
Metropolis--Hastings ratio, is the probability that a particular component of
a Gaussian vector is the largest. ``Pointing at $e_{j}$'' means
$\delta x_{j} > \delta x_{k}$ for every competitor $k$, i.e.\ every other
component falls below $\delta x_{j}$. Conditioning on the value of
$\delta x_{j}$ and using independence across $k$:
\begin{equation}
  a_{i\to j}
  \;=\; \mathbb{P}\!\left(\delta x_{j} > \delta x_{k}\ \ \forall k \neq i,j\right)
  \;=\; \int_{-\infty}^{\infty}\!\!\mathrm{d}\delta x_{j}\;
        f_{\mu_{j},\sigma_{j}}(\delta x_{j})
        \prod_{k \neq i,j}
        \int_{-\infty}^{\delta x_{j}}\!\!\mathrm{d}\delta x_{k}\;
        f_{\mu_{k},\sigma_{k}}(\delta x_{k}),
  \label{eq:aij-raw}
\end{equation}
where $f_{\mu,\sigma}$ is the Gaussian density. The inner integrals are
Gaussian cumulatives, so this is
$\int \mathrm{d}\delta x_{j}\, f_{\mu_{j},\sigma_{j}}(\delta x_{j})
 \prod_{k} \Phi_{\mu_{k},\sigma_{k}}(\delta x_{j})$.
Standardising the outer variable, $\delta x_{j} = \mu_{j} + \sigma_{j}u$:
\begin{equation}
  \boxed{\;
  a_{i\to j}
  \;=\; \int_{-\infty}^{\infty}\! \phi(u)
        \prod_{k \neq i,j}
        \Phi\!\left(\alpha_{kj} + \beta_{kj}\,u\right) \mathrm{d}u,
  \qquad
  \alpha_{kj} = \frac{\mu_{j}-\mu_{k}}{\sigma_{k}},
  \qquad
  \beta_{kj} = \frac{\sigma_{j}}{\sigma_{k}}.
  \;}
  \label{eq:aij}
\end{equation}
with $\phi$, $\Phi$ the standard normal density and CDF. Note that
$\alpha_{kj}$ is divided by $\sigma_{k}$ --- the standard deviation of the
competitor being compared against --- not by $\sigma_{j}$.

\paragraph{The specialisation used here.} In this sampler $\sigma$ is a
global scalar: $\delta x = \dt\,p/\Mass - \frac{\dt^{2}}{2\Mass}\gradU$ with
isotropic momentum $p\sim\mathcal{N}(0,\Mass T I)$, so every component has the
same variance $\dt^{2}T/\Mass$ and
\begin{equation}
  \beta_{kj} = \frac{\sigma_{j}}{\sigma_{k}} = 1 \quad \text{identically},
  \qquad
  \alpha_{kj} = \frac{\mu_{j}-\mu_{k}}{\sigma},
\end{equation}
which is what \path{utils/proposals.py} implements (it takes $\sigma$ as a
scalar argument). Dropping $\beta$ is therefore correct \emph{here} but not in
general: verified numerically against brute-force Monte Carlo,
\eqref{eq:aij} reproduces the true pointing probabilities to within Monte
Carlo error in both cases, whereas using a single $\sigma$ when the
$\sigma_{k}$ genuinely differ is wrong by up to $0.135$ in absolute
probability on a $K=6$ test. Any future variant with a per-class or per-site
step size reinstates $\beta$.

With constant $\sigma$ this is a Thurstone/probit discrete-choice
model---the probit analogue of the softmax-based proposal used by
Gibbs-With-Gradients. It is an exact probability, not merely something
proportional to one: summing \eqref{eq:aij} over all admissible $j$ at a fixed
site gives $1$, which is checked numerically in
\path{tests/test_proposals.py}.

\subsection{How it is evaluated}

\path{log_pointing_prob} evaluates \eqref{eq:aij} by Gauss--Hermite
quadrature entirely in log space, using $n_{\mathrm{quad}}=40$ nodes:
\begin{equation}
  \log a_{i\to j}
  \;=\; \log \sum_{q} \exp\!\left[
          \log w_{q} + \sum_{k \neq j} \log \Phi\!\left(\alpha_{kj} + u_{q}\right)
        \right],
\end{equation}
computed with \texttt{torch.special.log\_ndtr} and \texttt{logsumexp} so that
no intermediate underflows even when $\alpha$ is large. The current residue is
removed by \emph{slicing} it out of the tensor being maximised over, not by
adding a large negative penalty---a deliberate change from an earlier version,
since a finite penalty could in principle be approached by
$\dt^{2}$-scaled gradients at large $\dt$, whereas a sliced-out class is
structurally absent.

\subsection{The independence assumption is exact here, not an approximation}

Equation \eqref{eq:aij} treats the $K$ components of $\delta x$ as
independent, and the project has long recorded this as ``a deliberate, small,
controlled approximation'' --- the reasoning being that mean-centring the
momenta and the displacement induces the covariance $\sigma^{2}(I-J/K)$
rather than $\sigma^{2}I$, a pairwise correlation of $-1/(K-1)$. That
description is \textbf{wrong for the quantity actually being computed}, and
the code is more correct than its own comment claimed.

The sampler centres twice --- once inside \texttt{\_extract\_momenta}
($p \leftarrow p - \overline{p}$) and once on the displacement
($\delta x \leftarrow \delta x - \overline{\delta x}$). Both subtract the
\emph{same scalar from every component}, and $\arg\max$ is invariant under a
common shift. The realised proposal is therefore the $\arg\max$ of the
\emph{uncentred} displacement, whose components are genuinely independent
$\mathcal{N}(\mu_{k},\sigma^{2})$. The induced correlation never reaches the
functional being evaluated. Slicing to the 19 competitors afterwards does not
disturb this, since the same scalar was subtracted from those components too.

Verified by simulating the exact pipeline --- centred momentum, centred
displacement, $\arg\max$ over the competitor slice, $2\times10^{6}$ draws:
\begin{center}
\begin{tabular}{l r}
\toprule
$\max|\text{empirical} - \text{formula}(\sigma)|$ & $0.00033$ \\
$\max|\text{empirical} - \text{formula}(\sigma\sqrt{1-1/K})|$ & $0.00252$ \\
Monte Carlo standard error & $0.00026$ \\
\bottomrule
\end{tabular}
\end{center}
The formula with the \emph{uncentred} $\sigma$ agrees within Monte Carlo
error. The ``corrected'' marginal $\sigma\sqrt{1-1/K}$ --- the natural thing
to reach for once one notices the centring --- is wrong by about ten standard
errors. Drawing uncentred momenta directly reproduces the centred pipeline's
$\arg\max$ frequencies to $0.00036$, confirming the shift-invariance argument
rather than merely asserting it.

So $a_{i\to j}$ as computed is exact, up to quadrature error, which is
separately verified: Gauss--Hermite with $n_{\mathrm{quad}}=40$ converged to
$\sim\!6$ significant figures even in the steepest case tested. The
observation that the Metropolis--Hastings correction only needs the forward
and reverse probabilities \emph{actually used to propose} still holds, and
remains the reason that even a genuinely approximate $a$ would not bias the
target distribution --- but it is not needed to defend this one.

% ===========================================================================
\section{\texorpdfstring{Why informedness scales as $\dt/\sqrt{T}$}{Why informedness scales as dt/sqrt(T)}}
\label{sec:scaling}
% ===========================================================================

\subsection{The derivation}

This is the central result of the note, and it is three lines. Inspecting
\eqref{eq:aij}, the \emph{only} way $\dt$, $T$ and $\Mass$ enter $a_{i\to j}$
is through the standardised gaps $\alpha_{kj} = (\mu_{j}-\mu_{k})/\sigma$.
Substituting \eqref{eq:musigma}:
\begin{align}
  \alpha_{kj}
  &= \frac{\mu_{j}-\mu_{k}}{\sigma}
   = \frac{-\dfrac{\dt^{2}}{2\Mass}\left[(\gradU)_{j} - (\gradU)_{k}\right]}
          {\dt\sqrt{T/\Mass}}
  \notag\\[4pt]
  &= \frac{\dt}{2\sqrt{\Mass T}}\,
     \left[(\gradU)_{k} - (\gradU)_{j}\right]
  \;=\;\frac{\dt \, g_{kj}}{2\sqrt{\Mass T}},
  \qquad g_{kj} \equiv (\gradU)_{k} - (\gradU)_{j}.
  \label{eq:alpha-scaling}
\end{align}
The $\dt^{2}$ of the drift divides by the $\dt$ of the noise, leaving one
power of $\dt$; the temperature enters only through $\sigma$, as
$\sqrt{T}$. Hence:
\begin{equation}
  \boxed{\;
    \pmatch \text{ depends on } (\dt, T, \Mass, \gradU)
    \text{ only through } \frac{\dt}{\sqrt{\Mass T}}
    \text{ and the gradient gaps } g_{kj}.
  \;}
  \label{eq:scaling-law}
\end{equation}

\subsection{Three consequences}

\paragraph{1. An exact invariance.} Any two settings with the same
$\dt/\sqrt{\Mass T}$ have \emph{identical} informedness. With $\Mass=1$
throughout this project, holding $\pmatch$ fixed while changing temperature
requires
\begin{equation}
  \dt' \;=\; \dt \,\sqrt{\frac{T'}{T}}.
  \label{eq:rescale}
\end{equation}
This is an exact statement about the proposal, not an approximation or a
fitted relation.

\paragraph{2. Hotter sampling is less informed at fixed $\dt$.} Raising $T$
inflates $\sigma$ without touching $\mu$, so the same $\dt$ buys strictly less
informedness at higher temperature. There is a certain physical justice to
this---a hotter chain should explore more freely---but it means $\dt$ is
\emph{not} a transferable constant, and any quoted $\dt$ is meaningless
without the $T$ it was measured at.

\paragraph{3. The trap, which we fell into.}
\texttt{test\_informedness\_vs\_dt.py} hard-codes its own $T=2.0$ as a module
constant, so every $\pmatch$ it reports is a $T\!=\!2$ number. Reading
$\dt=63$ off that table and using it for a run at $T=107$---which is what the
first version of the queued command did---changes $\alpha$ by
$1/\sqrt{107/2} = 1/7.31$ and lands the proposal essentially back on the
uniform floor. In the equivalence of \eqref{eq:rescale},
\begin{equation}
  \dt=63 \text{ at } T=107
  \;\;\equiv\;\;
  \dt = \frac{63}{\sqrt{107/2}} = 8.6 \text{ at } T=2,
\end{equation}
and at $\dt\!\approx\!8.6$ the measured $T\!=\!2$ curve gives
$\pmatch \approx 0.06$ against a uniform baseline of $0.0526$. The correct
value for the intended working point is
\begin{equation}
  \dt = 63 \times \sqrt{107/2} \;=\; 63 \times 7.3144 \;\approx\; 460 .
\end{equation}
This is now recorded in a comment on \path{DTS_TO_SWEEP} so that the next
reader does not repeat it. For the record, \texttt{protein\_g}'s long-standing
$\dt=1300$ needed no rescaling only because its runs used $T=2.0$, the same
temperature the informedness script assumes---the coupling was invisible for
as long as only one temperature was ever used.

% ===========================================================================
\section{\texorpdfstring{Choosing $\dt$}{Choosing dt}}
% ===========================================================================

\subsection{\texorpdfstring{Why not simply make $\dt$ large?}{Why not simply make dt large?}}

By (B), large $\dt$ drives $\pmatch \to 1$, i.e.\ a deterministic greedy pick.
That is undesirable for three independent reasons.

\begin{enumerate}
  \item \textbf{The gradient's pick is not that good.} By (A), $\dt$ cannot
        improve it. Measured top-1 accuracy---how often $j^{\star}$ is the
        true energetic optimum among the 19 competitors, by brute force---is
        $2/10$ on \texttt{protein\_g} and $10/30$ on \texttt{zero\_polymer}
        measured away from the energy floor. So a fully greedy proposal would
        commit to a substitution that is wrong roughly two thirds of the time,
        with no mechanism left for exploring the alternatives.
  \item \textbf{A deterministic proposal cripples the
        Metropolis--Hastings correction.} As $\pmatch \to 1$ the proposal
        becomes a near-deterministic map, the reverse probability $a_{j\to i}$
        for the move that undoes it becomes very small, and the Hastings ratio
        $a_{j\to i}/a_{i\to j}$ swings wildly. The chain's mobility collapses.
  \item \textbf{Specific sites saturate at wildly different rates.} The
        measured spread across sites is large---at $\dt=1000$ on
        \texttt{zero\_polymer}, $\pmatch$ ranges from $0.40$ to $1.00$
        depending on the site. A $\dt$ chosen for the mean already pins some
        sites at $1.0$.
\end{enumerate}

Equally, $\dt$ too small gives a uniform proposal: correct, but it throws away
the gradient information the whole construction exists to exploit, and the
Hastings term becomes $\approx 0$ because $a_{i\to j} \approx a_{j\to i}
\approx 1/\Ncomp$.

\subsection{The working point}

The target is therefore deliberately \emph{partial} informedness---meaningfully
above uniform, nowhere near greedy. The convention this project settled on is
\begin{equation}
  \pmatch \approx 0.15,
\end{equation}
roughly $3\times$ the uniform baseline of $1/19 = 0.0526$ and far from
saturation. It leaves the gradient a real but non-binding say, keeps all 19
competitors reachable within a realistic move budget, and keeps both
$a_{i\to j}$ and $a_{j\to i}$ comfortably away from their extremes so the
Hastings correction stays well-conditioned.

\subsection{Measured curves}

Table~\ref{tab:pmatch} is the measured $\pmatch(\dt)$ for both references,
both at $T=2.0$, pooled over 10 brute-forced sites. Note these are
\emph{closed-form} evaluations of \eqref{eq:aij} at the real measured
gradients, not Monte Carlo estimates: once $\gradU$ is known at a site (one
forward/backward pass), sweeping arbitrarily many $\dt$ values costs nothing
further. This is why extending the sweep downward was free.

\begin{table}[htbp]
\centering
\begin{tabular}{r cc cc}
\toprule
 & \multicolumn{2}{c}{\texttt{protein\_g} ($L=56$)}
 & \multicolumn{2}{c}{\texttt{zero\_polymer} ($L=566$)} \\
\cmidrule(lr){2-3}\cmidrule(lr){4-5}
$\dt$ & mean $\pmatch$ & (min, max) & mean $\pmatch$ & (min, max) \\
\midrule
    10 & ---    & ---                & 0.0647 & (0.056, 0.090) \\
    20 & ---    & ---                & 0.0788 & (0.058, 0.137) \\
    50 & ---    & ---                & 0.1285 & (0.068, 0.313) \\
   100 & ---    & ---                & 0.2132 & (0.084, 0.562) \\
   200 & ---    & ---                & 0.3479 & (0.121, 0.818) \\
   300 & 0.0725 & (0.060, 0.099)     & 0.4436 & (0.160, 0.933) \\
   600 & 0.0954 & (0.068, 0.160)     & 0.6229 & (0.275, 0.999) \\
  1000 & 0.1280 & (0.080, 0.256)     & 0.7447 & (0.403, 1.000) \\
  1300 & 0.1524 & (0.089, 0.331)     & 0.7930 & (0.470, 1.000) \\
  2000 & 0.2047 & (0.111, 0.483)     & 0.8507 & (0.503, 1.000) \\
  4000 & 0.3197 & (0.178, 0.743)     & 0.9156 & (0.520, 1.000) \\
  8000 & 0.4669 & (0.300, 0.948)     & 0.9506 & (0.541, 1.000) \\
 16000 & 0.6266 & (0.391, 1.000)     & 0.9581 & (0.581, 1.000) \\
 32000 & 0.7691 & (0.472, 1.000)     & 0.9659 & (0.659, 1.000) \\
\bottomrule
\end{tabular}
\caption{Measured $\pmatch(\dt)$, both at $T=2.0$, $\Mass=1$, pooled over 10
sites. Uniform baseline $1/19 = 0.0526$. The \texttt{protein\_g} sweep was
never run below $\dt=300$; the \texttt{zero\_polymer} range was extended down
to $\dt=10$ precisely because the old range turned out to sit entirely above
its working point.}
\label{tab:pmatch}
\end{table}

The two curves are displaced by about a factor of $20$ in $\dt$: at a given
$\dt$, \texttt{zero\_polymer} is far more informed. By
\eqref{eq:alpha-scaling} with $T$ and $\Mass$ equal, that displacement is
entirely attributable to larger gradient gaps $g_{kj}$ on the longer sequence,
which is consistent with the independently measured $\approx\!18\times$ larger
gradient components there.

\begin{table}[htbp]
\centering
\begin{tabular}{c rr r}
\toprule
target $\pmatch$ & $\dt$ at $T=2$ & $\dt$ at $T=107$ & regime \\
\midrule
$\approx 0.0526$ & ---   & ---    & uniform floor, no gradient information \\
$0.10$           &   33  &   240  & weakly informed \\
$\mathbf{0.15}$  & $\mathbf{63}$ & $\mathbf{460}$ & \textbf{the working point} \\
$0.25$           &  127  &   931  & strongly informed \\
$0.50$           &  394  &  2885  & approaching greedy \\
$0.79$           & 1281  &  9372  & effectively greedy \\
\bottomrule
\end{tabular}
\caption{Working points for \texttt{zero\_polymer}, with the $T=107$ column
obtained from the $T=2$ column by the exact rescaling \eqref{eq:rescale},
factor $\sqrt{107/2}=7.3144$. The $T=2$ values are linear interpolations
within Table~\ref{tab:pmatch} and so carry interpolation error; the rescaling
itself is exact.}
\label{tab:working-points}
\end{table}

\subsection{Recommended practice}

\begin{enumerate}
  \item Never quote $\dt$ without the $T$ it belongs to.
  \item Changing $T$ requires rescaling $\dt$ by $\sqrt{T'/T}$,
        \eqref{eq:rescale}, or informedness silently moves.
  \item Re-measure Table~\ref{tab:pmatch} for a new reference sequence. The
        curve depends on the gradient gaps, which are a property of the
        sequence and the energy, and they moved by $20\times$ between these
        two references.
  \item Check the per-site (min, max) spread, not just the mean. A mean of
        $0.15$ can hide sites already saturated at $1.0$.
\end{enumerate}

% ===========================================================================
\section{\texorpdfstring{What $\dt=460$ vs.\ $\dt=63$ showed}{What dt=460 vs. dt=63 showed}}
\label{sec:460v63}
% ===========================================================================

\subsection{The experiment}

\texttt{validate\_boltzmann\_toy.py} makes the sampler's target distribution
exactly checkable. It frees exactly two sites (30 and 52), pins all 564 others
at the reference, and enumerates every one of the $20^{2}=400$ states by brute
force, giving the exact Boltzmann distribution $p_{\mathrm{exact}} \propto
\exp(-U/T)$. It then runs the real \texttt{\_step} logic restricted to those
two sites and compares the empirical visit histogram against
$p_{\mathrm{exact}}$.

Both runs used $T=107$ (chosen from the participation-ratio calibration: the
inherited $T=2.0$ puts $p=1.0000$ on a single state and the chain cannot
move), $\Mass=1$, 20\,000 moves, 4000 discarded as burn-in. They differ
\emph{only} in $\dt$, and by \S\ref{sec:scaling} that means they differ only
in informedness:
\begin{center}
\begin{tabular}{llc}
\toprule
run & $\dt$ & informedness at $T=107$ \\
\midrule
informed & 460 & $\pmatch \approx 0.15$, the intended working point \\
uniform  &  63 & $\pmatch \approx 0.06$, essentially the $1/19$ floor \\
\bottomrule
\end{tabular}
\end{center}
The $\dt=63$ run exists because it was the erroneous value in the first
version of the queued command. Running it alongside the corrected one turned
the mistake into a controlled comparison: the same sampler, same energy, same
temperature, with the gradient information effectively switched off.

\subsection{Results}

\begin{table}[htbp]
\centering
\begin{tabular}{l rr}
\toprule
 & $\dt=460$ (informed) & $\dt=63$ ($\approx$ uniform) \\
\midrule
acceptance rate                  & 0.3060   & 0.2893 \\
$\mathrm{KL}(p_{\mathrm{emp}} \,\|\, p_{\mathrm{exact}})$
                                 & 0.0316   & 0.0338 \\
total variation distance         & 0.0787   & 0.0848 \\
distinct states visited          & 224/400  & 223/400 \\
\midrule
mean $U$ under $p_{\mathrm{exact}}$   & $+185.20$ & $+185.20$ \\
mean $U$ under $p_{\mathrm{emp}}$     & $+185.63$ & $+187.63$ \\
deviation                             & $+0.44$ \;(0.24\%) & $+2.43$ \;(1.31\%) \\
\bottomrule
\end{tabular}
\caption{Boltzmann validation on \texttt{zero\_polymer}, $T=107$,
20\,000 moves, 16\,000 post-burn-in samples over 400 enumerated states.
$\log 400 = 5.99$ is the rough scale against which to read the KL.}
\label{tab:460v63}
\end{table}

\subsection{Reading the result}

\paragraph{Both pass; this validates the sampler at $L=566$.} The top states
come back in essentially the right order and the right proportions, and mean
$U$---the most robust single summary, being an average over the whole
distribution rather than a per-state frequency---agrees to $0.24\%$ at the
informed $\dt$. Neither run shows the decisive failure signature the script
warns about: mass piled on states the exact distribution calls unlikely, or
the true top states missed. This was the last outstanding safety check and the
only one that exercises accept/reject end to end.

\paragraph{The $\dt$ error cost coverage, not correctness.} This is the
substantive methodological lesson of the pair. A near-uniform proposal still
targets the same Boltzmann distribution---the Metropolis--Hastings correction
does not care how good the proposal is, only that its forward and reverse
probabilities are computed consistently. So the $\dt=63$ run converged
perfectly respectably. But with $a_{i\to j} \approx a_{j\to i} \approx
1/\Ncomp$, the Hastings term
\begin{equation}
  \log a_{j\to i} - \log a_{i\to j} \;\approx\; 0,
\end{equation}
which means that run never meaningfully exercised the reverse-probability
pass---the second backward pass at the candidate point, and the asymmetry
correction it feeds. That is the most intricate part of the sampler and the
part most likely to harbour a subtle bug, and it is the reason this test
exists. Had only the $\dt=63$ run been done, it would have been reported as a
pass while leaving that machinery untested. The informed run is the one that
actually tests it.

\paragraph{Informedness helps, but modestly, and not via acceptance.} The
informed run is better on every metric, but only slightly, and acceptance
barely moved ($0.306$ vs.\ $0.289$). Informedness is not acting as an
acceptance-rate lever here. The natural reading is that on a 400-state
problem with a $\approx\!41$-state effective support, a uniform proposal is
already a decent explorer, so there is little headroom; the informed
proposal's advantage should be expected to matter much more on the full
566-site problem, where the competitor set per move is the same 19 but the
state space is vastly larger. That expectation is untested.

\subsection{Is the residual discrepancy bias or noise?}
\label{sec:noise}

The KL values are small but not zero, and it is worth being precise about
whether that residue is bias. The participation ratio of
$p_{\mathrm{exact}}$ at $T=107$ is
\begin{equation}
  k_{\mathrm{eff}} = \Big(\textstyle\sum_{i} p_{i}^{2}\Big)^{-1} = 41.4
\end{equation}
effective states out of 400. For $n$ \emph{independent} samples from the exact
distribution, the expected KL of the empirical histogram is approximately
\begin{equation}
  \mathbb{E}\!\left[\mathrm{KL}\right] \approx \frac{k_{\mathrm{eff}}-1}{2n}
  = \frac{40.4}{2 \times 16000} = 0.00126,
\end{equation}
about $25\times$ smaller than observed. But MCMC samples are not
independent. Inverting the same expression for the effective sample size
actually implied:
\begin{equation}
  \mathrm{ESS} \approx \frac{k_{\mathrm{eff}}-1}{2\,\mathrm{KL}}
  = \begin{cases}
      639 & (\dt=460)\\
      598 & (\dt=63)
    \end{cases}
  \qquad\Longrightarrow\qquad
  \tau_{\mathrm{auto}} = \frac{16000}{\mathrm{ESS}} \approx 25\text{--}27 \text{ moves}.
\end{equation}
An autocorrelation time of $\approx\!25$ moves is entirely ordinary for a
two-site chain at $30\%$ acceptance: each site is re-proposed roughly every
two moves and the proposal is accepted about $30\%$ of the time, so the state
decorrelates over tens of moves. Both KL values are therefore what a
\emph{correct} sampler limited by chain autocorrelation looks like. This is an
estimate, not a measurement---$\tau_{\mathrm{auto}}$ was inferred from the KL
rather than computed from the chain.

\paragraph{The falsifiable check, and its outcome.} If the sampler is
unbiased, the KL is pure finite-sample noise and must fall as $1/n$, while the
total variation distance---a different functional of the same histogram---must
fall as $1/\sqrt{n}$. Predicting forward from the $n=16000$ run gave
$\mathrm{KL}(32000) \approx 0.016$; a KL that instead plateaued near $0.03$
would have indicated genuine bias. The run was done
(\path{zero_polymer/validate_boltzmann_results_dt460_M40000.txt}, 40\,000
moves, 32\,000 post burn-in) and \textbf{both laws hold}:

\begin{table}[htbp]
\centering
\begin{tabular}{l rr r r l}
\toprule
 & $n=16000$ & $n=32000$ & ratio & predicted & \\
\midrule
$\mathrm{KL}(p_{\mathrm{emp}}\,\|\,p_{\mathrm{exact}})$
  & 0.0316 & 0.0169 & 1.87 & 2.00 \;($1/n$)        & pass \\
total variation distance
  & 0.0787 & 0.0569 & 1.38 & 1.41 \;($1/\sqrt{n}$) & pass \\
\midrule
$\tau_{\mathrm{auto}}$ (inferred)
  & 25.0 & 26.8 & --- & run-length independent & pass \\
acceptance rate
  & 0.3060 & 0.3092 & --- & unchanged & pass \\
\bottomrule
\end{tabular}
\caption{The $1/n$ bias check. KL and total variation are \emph{independent}
functionals of the same histogram with different convergence rates, so
matching both is considerably stronger evidence than matching either alone: a
biased sampler would have each converge to its own nonzero constant rather
than keep falling on schedule. The inferred autocorrelation time is a property
of the chain and must not depend on run length---it does not.}
\label{tab:bias-check}
\end{table}

\noindent
This settles the question to the resolution available: the residual
discrepancy is finite-sample noise, behaving quantitatively as finite-sample
noise must, and there is no evidence of bias. The claim can now be stated more
strongly than \emph{no detectable bias at this resolution}---though note that
no finite experiment can establish exact unbiasedness, and what has been shown
is that the error decays on the schedule an unbiased sampler requires.

\paragraph{One number that did not improve, and why that is fine.} Mean $U$
moved \emph{away} from the exact value between the two runs ($+0.44 \to
+1.70$). This is not a regression. Mean $U$ is a single scalar statistic
carrying its own Monte Carlo error, not a convergence metric: with
$\mathrm{ESS}\approx 1200$ and $\mathrm{sd}_{p}(U)$ of order 150--200, its
standard error is roughly $5$, so both deviations sit comfortably inside one
standard error. The $0.24\%$ agreement at $n=16000$ reported in
Table~\ref{tab:460v63} was therefore partly luck, and should not be read as a
precision claim; the KL and total variation, which aggregate over the whole
distribution, are the trustworthy convergence diagnostics here.

% ===========================================================================
\section{Caveats and open items}
% ===========================================================================

\begin{enumerate}
  \item \textbf{The $T$ window is from a two-site subproblem.} The usable
        range $T \approx 32$--$322$, and hence the $\dt \approx 460$ derived
        from it, were obtained with 564 of 566 residues pinned at the
        reference. A full-length chain explores a far wider energy range, so
        that window is a starting bracket, not a final answer.
        \texttt{zero\_polymer} needs its own temperature sweep.
  \item \textbf{Informedness curves are floor measurements.} The gradients
        behind Table~\ref{tab:pmatch} were computed at a $\mathrm{seq\_A}$
        only 5 mutations from the reference---$0.9\%$ drift at $L=566$, where
        the energy sits essentially at its floor. Measured at a comparable
        drift to \texttt{protein\_g} ($\approx\!8.8\%$), the gradient's
        behaviour changes substantially: top-1 rises from $2/10$ to $10/30$
        and the mean true $\Delta U$ at the gradient's pick goes from $+474$
        to $-360$. Gradient \emph{gaps}, and hence the whole $\pmatch(\dt)$
        curve, should be re-measured in the bulk before
        Table~\ref{tab:working-points} is trusted for a production run.
  \item \textbf{The independence approximation} of \S2 is the one inexactness
        in $a_{i\to j}$. Unquantified against the exact orthant integral.
  \item \textbf{Per-site saturation} is wide (Table~\ref{tab:pmatch}), so a
        single global $\dt$ is necessarily a compromise. Whether a per-site
        $\dt$ would help, and whether it could be made to preserve detailed
        balance, is unexplored.
  \item \textbf{Resolved:} the $1/n$ scaling of the KL (\S\ref{sec:noise})
        was the one outstanding validation item at the time of writing. It has
        since been run and passed, together with the $1/\sqrt{n}$ scaling of
        the total variation distance---see Table~\ref{tab:bias-check}. What
        remains open on the validation side is only that this is a two-site
        subproblem: the full 566-site chain has never been validated against
        anything exact, and cannot be, since its state space is not
        enumerable.
\end{enumerate}

\end{document}
